%%%%%%%%%%%%%%%%%%%%%%%%%%%%%%%%%%%%%%%%%%%%%%%%%%%%%%%%%%%%%%
%% Chapter 6: Sentence-level Translation (video2text)
%%%%%%%%%%%%%%%%%%%%%%%%%%%%%%%%%%%%%%%%%%%%%%%%%%%%%%%%%%%%%%

\section{Sentence-level Translation (video2text)}
\label{chap:sentence-level-translation}

Chapters~\ref{chap:isolated-gloss-recognition} and~\ref{chap:text-translation}
address, respectively, isolated video-to-gloss recognition and gloss-text
translation, each operating on a single, pre-segmented gloss or a short gloss
sequence. Practical sign language translation, however, requires processing
continuous, unsegmented signing sequences in which gloss boundaries are unknown
and must be inferred jointly with gloss identity. This chapter describes the
end-to-end pipeline that integrates the two preceding modules into a full
video-to-text system.

\subsection{End-to-end pipeline}
\label{sec:sentence-level-pipeline}

We plan to extend the isolated recognition system of
Chapter~\ref{chap:isolated-gloss-recognition} toward continuous recognition
using Connectionist Temporal Classification (CTC), a sequence-to-sequence
training objective that marginalises over all valid alignments between an input
sequence and an output label sequence of unknown alignment. CTC introduces a
blank token to model inter-gloss transitions and trains the encoder to emit
gloss labels at the appropriate temporal positions without requiring
frame-level annotation. This extension would transform the current isolated
recognition architecture into a continuous sign language recognition (CSLR)
pipeline, significantly expanding its practical applicability while retaining
the learned embedding space of Chapter~\ref{chap:isolated-gloss-recognition} as
its representational backbone.

\todo{PLACEHOLDER: Describe the concrete end-to-end pipeline architecture once
implemented: how continuous video is segmented or processed frame-by-frame,
how the CTC-based (or otherwise) continuous recognition stage produces a gloss
sequence, and how that sequence is handed off to the gloss-to-text stage of
Chapter~\ref{chap:text-translation}.}

\subsection{Integration of modules from Chapters~\ref{chap:isolated-gloss-recognition} and~\ref{chap:text-translation}}
\label{sec:sentence-level-integration}

\todo{PLACEHOLDER: Describe how the isolated gloss recognition model
(Chapter~\ref{chap:isolated-gloss-recognition}) and the gloss-text translation
model (Chapter~\ref{chap:text-translation}) are composed into a single
inference pipeline -- interface between the two modules, handling of
recognition uncertainty/ambiguity when passed downstream to translation, and
any joint fine-tuning or error-correction mechanisms between the two stages.}

\subsection{Experiments and results}
\label{sec:sentence-level-experiments}

\todo{PLACEHOLDER: Report end-to-end evaluation of the full video-to-text
pipeline on continuous, sentence-level sign language video, including the
dataset used, end-to-end metrics (e.g.\ Word Error Rate for the recognition
stage, BLEU for the translation stage, and any combined video-to-text metric),
and an error analysis distinguishing recognition errors from translation
errors.}
